%%%%%%%%%%%%%%%%%%%%%%%%%%%%%%%%%%%%%%%%%%%%%%%%%%%%%%%%%%%%%%%%%%%%%%%%%%%%%%%%%%%%%%
\section{Model}
\label{sec:perliminaries}
%%%%%%%%%%%%%%%%%%%%%%%%%%%%%%%%%%%%%%%%%%%%%%%%%%%%%%%%%%%%%%%%%%%%%%%%%%%%%%%%%%%%%%
%
%%=====================================================================================
%\subsection{Model}
%\label{sec:model}
%%=====================================================================================
We assume a set of \numprocesses nodes, such that each one of them represents an independent trust authority.
A computationally bound adversary controls the network and may corrupt up to \numfaults nodes, such that $\numprocesses \ge 3\numfaults + 1$.
We refer to the corrupted nodes as Byzantine, as they may fail arbitrarily, and to the non-corrupted nodes as correct. 

Each node participates in various types of roles/activities, encapsulated by processes hosted on the same node:
\begin{itemize}
	\item \textit{Validation}: Validator processes are responsible for populating the log with blocks. Each node hosts a single validator process.
	\item \textit{Work}: Worker processes are responsible for transaction dissemination. Each node hosts at least one worker process.
	\item \textit{Ticketing}: A ticketmaster process is a functionality responsible for assigning log slots to validators (see \Cref{sec:ticketmaster} for dertails). A ticketmaster may be implemented by process hosted by nodes, in which case each node hosts at most one ticketmaster process. 
\end{itemize}
Below, we will interchangably refer to a node and to the processes it hosts, since there is a 1-1 mapping between them.
Furthermore, we assume that the adversary controls any process of a corrupted node.

Unless mentioned otherwise, processes communicate through authenticated point-to-point channels.
We assume an eventually synchronous network~\cite{DworkLS88}
such that the communication between any pair of correct processes is asynchronous before a global stabilization time ($GST$), unknown to the processes, 
when the communication becomes synchronous.
Finally, we assume a public key infrastructure (PKI) such that each process $i$ is equipped with a public-private key pair $PK_i$,$SK_i$ used for authentication.
We use $\langle \mathit{[field]}\ldots, \sigma_i\rangle$ to denote messages carrying a signature $\sigma_i$ by $i$.